\documentclass[12pt,a4paper]{article}
\usepackage[margin=2.5cm]{geometry}
\usepackage{graphicx,hyperref,listings,xcolor,booktabs,tabularx,fancyhdr,titlesec,enumitem,amsmath}

\hypersetup{colorlinks=true,linkcolor=blue,urlcolor=cyan}
\definecolor{codebg}{RGB}{13,21,38}
\definecolor{codetext}{RGB}{200,210,230}
\definecolor{accent}{RGB}{0,180,220}

\lstset{
  backgroundcolor=\color{codebg},
  basicstyle=\ttfamily\small\color{codetext},
  keywordstyle=\color{accent}\bfseries,
  commentstyle=\color{gray},
  stringstyle=\color{green!60!white},
  breaklines=true, frame=single, rulecolor=\color{gray!40},
  numbers=left, numberstyle=\tiny\color{gray}, numbersep=8pt
}

\pagestyle{fancy}
\fancyhf{}
\rhead{\textcolor{gray}{\small Mini Project \#13 — Weather Station}}
\lhead{\textcolor{gray}{\small Software Tools Lab}}
\cfoot{\thepage}

\title{
  \vspace{-1cm}
  {\Large \textbf{Mini Project \#13}}\\[0.4cm]
  {\LARGE \textbf{Real-Time Weather Station Dashboard}}\\[0.3cm]
  {\large A Full-Stack Web Application Covering Experiments 1--12}\\[0.5cm]
  \hrule
}
\author{Software Tools Laboratory}
\date{\today}

\begin{document}
\maketitle\thispagestyle{fancy}
\tableofcontents\newpage

%% ─────────────────────────────────────────────
\section{Project Overview}
This mini project implements a \textbf{Real-Time Weather Station Dashboard} — a full-stack application that synthesises concepts from all twelve laboratory experiments. A Python backend simulates atmospheric sensor readings and stores them in a SQLite database. A Flask API serves the data as both JSON and XML. A JavaScript-powered web dashboard polls the API every three seconds via AJAX and renders live charts and gauges. A Tkinter desktop GUI allows manual submission of readings. The entire stack is containerised with Docker.

\subsection{Objectives}
\begin{enumerate}[label=\arabic*.]
  \item Model the system using Data Flow, Entity-Relationship, and UML diagrams.
  \item Implement forward and reverse engineering via database schema generation and extraction.
  \item Automate server startup using a Linux Bash shell script.
  \item Check network interface availability with a Python script.
  \item Build a dynamic JavaScript web dashboard with AJAX data polling.
  \item Implement Flask web forms with client- and server-side validation.
  \item Deploy the application using Docker and Docker Compose.
  \item Develop an AJAX-enabled Rich Internet Application with JSON and XML data formats.
  \item Create a Python Tkinter GUI for manual sensor data entry.
  \item Document the project using \LaTeX.
\end{enumerate}

%% ─────────────────────────────────────────────
\section{Experiment 1: Data Flow Diagram and ER Diagram}

\subsection{Data Flow Diagram (DFD)}
\textbf{Level 0 — Context Diagram:}
The system has three external entities: \textit{Sensor Simulator}, \textit{Web Browser}, and \textit{Tkinter GUI}. All data flows through the central process: \textit{Weather Station System}.

\textbf{Level 1 DFD — Major Processes:}
\begin{enumerate}
  \item \textbf{P1 — Generate Reading:} Sensor Simulator $\rightarrow$ generates reading $\rightarrow$ SQLite DB
  \item \textbf{P2 — Serve API:} SQLite DB $\rightarrow$ Flask reads $\rightarrow$ JSON/XML response $\rightarrow$ Browser
  \item \textbf{P3 — Submit Manual:} Tkinter GUI $\rightarrow$ POST /api/submit $\rightarrow$ SQLite DB
  \item \textbf{P4 — Generate Alert:} Reading crosses threshold $\rightarrow$ Alerts table $\rightarrow$ Dashboard
  \item \textbf{P5 — Register Station:} Browser form $\rightarrow$ POST /register $\rightarrow$ Stations table
\end{enumerate}

\subsection{Entity-Relationship Diagram}
Three entities with the following attributes and relationships:

\begin{center}
\begin{tabular}{lll}
\toprule
\textbf{Entity} & \textbf{Attributes} & \textbf{Key} \\
\midrule
Station & id, name, location, latitude, longitude, registered\_at & id (PK) \\
Reading & id, station\_id, temperature, humidity, pressure, wind\_speed, wind\_direction, recorded\_at, source & id (PK) \\
Alert   & id, station\_id, reading\_id, alert\_type, message, triggered\_at & id (PK) \\
\bottomrule
\end{tabular}
\end{center}

\textbf{Relationships:}
\begin{itemize}
  \item \textbf{Station} \textit{generates} \textbf{Reading} — One-to-Many (1:N)
  \item \textbf{Reading} \textit{triggers} \textbf{Alert} — One-to-Zero-or-One (1:0..1)
\end{itemize}

%% ─────────────────────────────────────────────
\section{Experiment 2: UML Diagrams}

\subsection{Use Case Diagram}
\textbf{Actors:} Web User, GUI Operator, System (auto-sensor).

\textbf{Use Cases:}
\begin{itemize}
  \item Web User: View Dashboard, Toggle JSON/XML, Register Station, View Alerts
  \item GUI Operator: Launch GUI, Submit Manual Reading
  \item System: Auto-Generate Reading, Check Thresholds, Insert Alert
\end{itemize}

\subsection{Class Diagram}
\begin{lstlisting}[language=Python, caption=Class Structure]
class SensorSimulator:
    state: dict
    + generate_reading() -> dict
    + start_background_sensor(interval: int)

class Database:
    DB_PATH: str
    + init_db()
    + get_connection() -> Connection
    + reverse_engineer_schema() -> dict

class FlaskApp:
    + index() -> HTML
    + api_current() -> JSON
    + api_readings() -> JSON
    + api_readings_xml() -> XML
    + api_submit() -> JSON
    + register_station() -> HTML

class StationGUI:
    root: Tk
    + _build_form()
    + _submit()
    + _fetch_stations()
\end{lstlisting}

\subsection{Sequence Diagram — AJAX Poll}
\begin{enumerate}
  \item Browser JS calls \texttt{fetch('/api/current')} every 3s
  \item Flask queries SQLite: \texttt{SELECT * FROM readings ORDER BY id DESC LIMIT 1}
  \item Flask returns JSON response
  \item \texttt{app.js} calls \texttt{updateCards(data)} and \texttt{updateChart(data)}
  \item Chart.js re-renders; DOM values update with animation
\end{enumerate}

%% ─────────────────────────────────────────────
\section{Experiment 3: Forward and Reverse Engineering}

\subsection{Forward Engineering}
The ER diagram is translated into SQLite DDL statements in \texttt{server/database.py}:

\begin{lstlisting}[language=SQL, caption=Forward Engineering — DDL from ER]
CREATE TABLE IF NOT EXISTS stations (
    id           INTEGER PRIMARY KEY AUTOINCREMENT,
    name         TEXT NOT NULL UNIQUE,
    location     TEXT NOT NULL,
    latitude     REAL,
    longitude    REAL,
    registered_at TEXT DEFAULT (datetime('now'))
);
CREATE TABLE IF NOT EXISTS readings (
    id             INTEGER PRIMARY KEY AUTOINCREMENT,
    station_id     INTEGER NOT NULL,
    temperature    REAL NOT NULL,
    humidity       REAL NOT NULL,
    pressure       REAL NOT NULL,
    wind_speed     REAL NOT NULL,
    wind_direction TEXT DEFAULT 'N',
    recorded_at    TEXT DEFAULT (datetime('now')),
    source         TEXT DEFAULT 'auto',
    FOREIGN KEY (station_id) REFERENCES stations(id)
);
\end{lstlisting}

\subsection{Reverse Engineering}
The \texttt{/api/schema} endpoint extracts table metadata using \texttt{PRAGMA table\_info()} and returns a JSON structure representing the inferred ER model — simulating reverse engineering of the database.

%% ─────────────────────────────────────────────
\section{Experiment 4: Linux Shell Programming}

The file \texttt{scripts/start.sh} automates the full startup process:

\begin{lstlisting}[language=bash, caption=Shell Script Highlights]
#!/usr/bin/env bash
# 1. Check Python installation
command -v python3 || { echo "Python not found"; exit 1; }
# 2. Create and activate virtual environment
python3 -m venv .venv && source .venv/bin/activate
# 3. Install dependencies
pip install -q -r requirements.txt
# 4. Run network check (Experiment 5)
python3 server/network_check.py
# 5. Initialise database
python3 -c "from server.database import init_db; init_db()"
# 6. Start Flask in background, tail logs
nohup python3 server/app.py > server.log 2>&1 &
tail -f server.log
\end{lstlisting}

%% ─────────────────────────────────────────────
\section{Experiment 5: Python Network Interface Checker}

\texttt{server/network\_check.py} performs three checks:

\begin{enumerate}
  \item \textbf{Interface Listing} — uses \texttt{psutil.net\_if\_addrs()} to list all IPv4 interfaces with up/down status.
  \item \textbf{Internet Check} — opens a socket to Google DNS (\texttt{8.8.8.8:53}).
  \item \textbf{Server Reachability} — attempts \texttt{socket.create\_connection(('127.0.0.1', 5000))}.
\end{enumerate}

\begin{lstlisting}[language=Python, caption=Socket-based server reachability check]
import socket
def check_server(host='127.0.0.1', port=5000, timeout=3):
    try:
        sock = socket.create_connection((host, port), timeout=timeout)
        sock.close()
        print(f"Server REACHABLE at {host}:{port}")
        return True
    except ConnectionRefusedError:
        print("Server is NOT running")
        return False
\end{lstlisting}

%% ─────────────────────────────────────────────
\section{Experiment 6: JavaScript Dynamic Web Pages}

\texttt{static/app.js} implements a fully dynamic single-page dashboard:

\begin{itemize}
  \item \textbf{AJAX Polling:} \texttt{setInterval(poll, 3000)} — calls \texttt{fetch('/api/current')} every 3 seconds.
  \item \textbf{DOM Manipulation:} Card values, progress bars, and the compass needle update on every poll.
  \item \textbf{Chart.js:} A rolling 30-point line chart renders temperature, humidity, and pressure history.
  \item \textbf{JSON/XML Toggle:} Users switch between data sources; the raw API response displays in a code panel.
  \item \textbf{XML Parsing:} \texttt{DOMParser} parses the XML feed client-side and extracts values identically to JSON mode.
\end{itemize}

%% ─────────────────────────────────────────────
\section{Experiment 7: Web Programming — Form Design and Validation}

The \texttt{/register} route demonstrates both client- and server-side validation:

\textbf{Client-side (JavaScript):}
\begin{lstlisting}[language=JavaScript, caption=Client-side validation]
form.addEventListener('submit', function(e) {
    if (name.length < 3) { alert('Name too short'); e.preventDefault(); }
    if (isNaN(lat) || lat < -90 || lat > 90) {
        alert('Invalid latitude'); e.preventDefault();
    }
});
\end{lstlisting}

\textbf{Server-side (Flask/Python):}
\begin{lstlisting}[language=Python, caption=Server-side validation in Flask]
if not name:
    errors['name'] = 'Station name is required.'
elif len(name) < 3:
    errors['name'] = 'Name must be at least 3 characters.'
try:
    lat = float(lat_str)
    if not -90 <= lat <= 90:
        errors['latitude'] = 'Latitude must be -90 to 90.'
except ValueError:
    errors['latitude'] = 'Enter a valid decimal number.'
\end{lstlisting}

%% ─────────────────────────────────────────────
\section{Experiment 8: Docker Deployment}

\subsection{Dockerfile}
\begin{lstlisting}[caption=Dockerfile]
FROM python:3.11-slim
WORKDIR /app
COPY requirements.txt .
RUN pip install --no-cache-dir -r requirements.txt
COPY server/ ./server/
COPY templates/ ./templates/
COPY static/ ./static/
EXPOSE 5000
CMD ["python", "server/app.py"]
\end{lstlisting}

\subsection{Docker Compose}
\begin{lstlisting}[caption=docker-compose.yml]
version: '3.9'
services:
  weather-server:
    build: .
    ports:
      - "5000:5000"
    volumes:
      - ./data:/app/data
    restart: unless-stopped
    healthcheck:
      test: ["CMD","python","-c",
             "import urllib.request; urllib.request.urlopen('http://localhost:5000/')"]
      interval: 30s
      retries: 3
\end{lstlisting}

%% ─────────────────────────────────────────────
\section{Experiment 9: AJAX-Enabled RIA with JSON and XML}

The application provides dual data format support:

\textbf{JSON endpoint:} \texttt{GET /api/readings} returns:
\begin{lstlisting}[language=Python, caption=Sample JSON response]
{
  "count": 20,
  "readings": [
    {"id": 42, "temperature": 28.5, "humidity": 61.2,
     "pressure": 1013.1, "wind_speed": 14.3,
     "wind_direction": "NE", "recorded_at": "2024-01-15 10:23:05"}
  ]
}
\end{lstlisting}

\textbf{XML endpoint:} \texttt{GET /api/readings.xml} returns:
\begin{lstlisting}[language=XML, caption=Sample XML response]
<?xml version="1.0" ?>
<WeatherData generated_at="2024-01-15T10:23:05Z">
  <Reading>
    <temperature>28.5</temperature>
    <humidity>61.2</humidity>
    <pressure>1013.1</pressure>
    <wind_speed>14.3</wind_speed>
    <wind_direction>NE</wind_direction>
  </Reading>
</WeatherData>
\end{lstlisting}

%% ─────────────────────────────────────────────
\section{Experiment 10: Python GUI — Tkinter}

\texttt{gui/station\_gui.py} provides a desktop application with:
\begin{itemize}
  \item Input fields for Temperature, Humidity, Pressure, Wind Speed, Wind Direction
  \item Station selector (populated via \texttt{GET /api/stations})
  \item Client-side range validation before submission
  \item \texttt{urllib.request} POST to \texttt{/api/submit}
  \item Status bar showing success/failure with colour coding
\end{itemize}

%% ─────────────────────────────────────────────
\section{Experiment 11: Document Creation Using \LaTeX}
This document itself is the \LaTeX\ deliverable for Experiment 11. It was authored in \texttt{docs/report.tex} and compiled via Overleaf. It covers the complete project documentation including diagrams, code listings, and analysis.

%% ─────────────────────────────────────────────
\section{Experiment 12: Software Management Tools}

\subsection{Version Control (Git)}
\begin{lstlisting}[language=bash, caption=Git workflow]
git init
git add .
git commit -m "feat: initial weather station project"
# Feature branch workflow
git checkout -b feature/xml-endpoint
git commit -m "feat: add XML API endpoint"
git checkout main && git merge feature/xml-endpoint
\end{lstlisting}

\subsection{Bug Reporting (Bugzilla-style)}
A sample bug report for this project:

\begin{center}
\begin{tabularx}{\textwidth}{lX}
\toprule
\textbf{Field} & \textbf{Value} \\
\midrule
Bug ID     & WS-001 \\
Summary    & Compass needle resets to North on XML mode switch \\
Severity   & Minor \\
Priority   & P3 \\
Component  & Frontend / app.js \\
Steps      & 1. Switch to XML mode. 2. Observe compass direction resets. \\
Expected   & Compass should retain last known direction \\
Actual     & Compass needle jumps to 0° (North) \\
Fix        & Cache last wind\_direction; use cached value if XML parse returns null \\
Status     & Resolved \\
\bottomrule
\end{tabularx}
\end{center}

%% ─────────────────────────────────────────────
\section{Testing and Results}

\begin{center}
\begin{tabular}{llll}
\toprule
\textbf{Test} & \textbf{Input} & \textbf{Expected} & \textbf{Result} \\
\midrule
Server Start & python server/app.py & HTTP 200 on port 5000 & \textcolor{green!60!black}{Pass} \\
JSON API & GET /api/current & JSON with all fields & \textcolor{green!60!black}{Pass} \\
XML API & GET /api/readings.xml & Valid XML document & \textcolor{green!60!black}{Pass} \\
Form Validation & Empty station name & Error message shown & \textcolor{green!60!black}{Pass} \\
GUI Submit & Valid reading via Tkinter & Reading appears on dashboard & \textcolor{green!60!black}{Pass} \\
Docker Deploy & docker-compose up & Dashboard at localhost:5000 & \textcolor{green!60!black}{Pass} \\
Network Check & python network\_check.py & Interface list printed & \textcolor{green!60!black}{Pass} \\
Alert Trigger & temperature > 42°C & Alert banner shown & \textcolor{green!60!black}{Pass} \\
\bottomrule
\end{tabular}
\end{center}

%% ─────────────────────────────────────────────
\section{Conclusion}
The Real-Time Weather Station Dashboard successfully integrates all twelve laboratory experiments into a cohesive, functional application. The project demonstrates a complete software development lifecycle: from design (DFD, ER, UML) through implementation (Python, Flask, JavaScript, Tkinter, Shell) to deployment (Docker) and documentation (\LaTeX). The live AJAX dashboard, dual JSON/XML data formats, and containerised deployment make it a practical example of modern full-stack development.

\end{document}
